\section{Problem 1}
\textit{Consider a causal system given by the following system function}
\[H(z) = 1 + 3.5z^{-1}+1.5z^{-2}\]
\textit{And assume that the causal systtem is used as a model for a communcation channel that will
distort signals passing through it.}

\subsection{1.}
\textit{Design an equalizer that can mitigate the distortion}
The system function is being described as causal, meaning that it cannot depend on future
inputs, only on previous inputs $x[n], x[n-1], x[n-2], ...$. In non-minimum-phase systems, the
distortion from systems $H(z)$ can be \textit{undone} by cascading an equalizer, or filter, with 
the system function, such that when multiplied by its inverse, its effects will be undone.
$H(z)H_{inv}(z) = 1$. %side 254 i bogen
From the given system function, the zeros and poles can be extracted. A stable and causal 
system, should have poles inside the unit circle, but the zeros does not need to be.
To find out what kind of equalizer we should construct, we can firstly analyze our system to
find these.

\[H(z) = 1 + 3.5z^{-1}+1.5z^{-2} \Rightarrow 
H(z)z^2 = z^2 + 3.5z^1 + 1.5 \Rightarrow
H(z) = \frac{z^2 + 3.5z^1 + 1.5}{z^2}\]\\
rewritten as \[\frac{(z+3)(z+0.5)}{z^2}\]\\
yielding zeroes in: $z=-3$,$z=-0.5$.
Here one zero lies within the unit circle, while the other does not.
This is not a problem for the original system.

Our equalizer should ideally be an inverse of the original system, such that 
\[H_{inv} = \frac{1}{H(z)} = \frac{A(z)}{B(z)}\] such that \[H(z) = \frac{B(z)}{A(z)}\]
however because of this, the zeroes from the original becomes the poles, which now resides outside
the unit circle making it unstable. Therefore we must decompose this system into a minimum-phase system
and an allpass system. 
\[H(z) = H_{min}(z)H_{ap}(z)\]
taking the deduced system function from earlier, this can be rewritten into an expression, 
that can be made into a system function that will be composed of our minimum-phase and all-pass systems
\[H(z) = \frac{(z+3)(z+0.5)}{z^2} \Rightarrow
\frac{z+0.5}{z} \frac{z+3}{z} \Rightarrow
(1+0.5z^{-1})(1+3z^{-1})\]
In this case, we have a zero $z=-3$ outside the unit circuit, and $z =\frac{1}{a}, |a|<1 \Rightarrow a = -\frac{1}{3}$
The composition ca be written as:
\[
H(z) = H_1(z)
\underbrace{(1-az^{-1})}_{\text{minimum-phase factor}}
\underbrace{
\frac{z^{-1}-a^*}{1-az^{-1}}
}_{\text{All-pass factor}}
\]
\[\Downarrow\]
\[
H(z) = 3(1+0.5z^{-1})
(1+\frac{1}{3}z^{-1})
\frac{z^{-1}+\frac{1}{3}}{1+\frac{1}{3}z^{-1}}
\]
Where $H_1(z] = (1+0.5z^{-1})$ and the additional factor of 3, added to the front of $H_1(z)$ 
is the scale factor, which the magnitude of a.

cleaning up the expressiong yields:
\[H(z)= 3(1+0.5z^{-1})(z^{-1}+\frac{1}{3})\]
All-pass filters affect phase, but not magnitude, therefore we can construct our filter assume
\[H_{eq=\frac{1}{H_{min}}} \Rightarrow \frac{1}{3(1+0.5z^{-1})
(1+\frac{1}{3}z^{-1})}\]
%husk scale factor på 3 a = 1/3, 

% OUTPUT = H(Z) * INPUT -> OUTPUT/INPUT = H(Z)
% ZEROS (OUTPUT)    - IF THIS GOES TO -> 0, THE ENTIRE THIS SHUTS DOWN
% ______________
% POLES (INPUT)     - IF THIS GOES TO -> 0, THE ENTIRE THINGS BLOWS UP TO INFINITY 